\section*{Chapter \ref{ch:rs:live-experiments} --- Live Experiments}

\begin{solution}{exo:rs:live-experiments:1}
$-0.55 \pm 1.96 \times 0.48$: from $-1.49$ to 0.39. The interval contains zero and savings three times larger than any plausible one: the experiment is
uninformative, not negative. It needed more orders, a larger share, or \omterm{def:rs:live-experiments:cuped}{CUPED}.
\end{solution}

\begin{solution}{exo:rs:live-experiments:2}
$z = (2\,485 - 2\,400)/\sqrt{24\,000 \times 0.1 \times 0.9} = 85/46.5 = 1.83$, a two-sided p-value of 0.07: no mismatch at 5\%. At 0.0001, stop reading the
estimate and find the cause (orders dropped, duplicated or re-routed in one arm, a filter applied after assignment, a server running an old assignment rule);
the comparison is not valid until the split is explained.
\end{solution}

\begin{solution}{exo:rs:live-experiments:3}
$2.80 \times 0.085 = 0.24$ basis point; the power against 0.3 is 0.94.
\end{solution}

\begin{solution}{exo:rs:live-experiments:4}
Allowed: the pre-trade estimate (fixed before assignment), the market's move over the scheduled window (outside the algorithm's reach), the stock's
previous-day shortfall (before assignment). Not allowed: the realised duration and the fill rate, which the algorithm changes; adjusting for them removes part
of the effect. With $R^2 = 0.51$ the variance is multiplied by 0.49: $77 \times 0.49 = 38$ days.
\end{solution}

\begin{solution}{exo:rs:live-experiments:5}
$n = 3^2 \times 23.1^2/(5^2(1-s))$: 194 orders and 8.1 days at 1\%, 214 orders and 0.89 day (348 minutes of a 6.5-hour session) at 10\%, 385 orders and
0.32 day at 50\%. Knight's defect sent child orders without regard to the parent's filled quantity: an invariant (child quantity never exceeds the parent's)
checked on each order, a position limit sized to the stage, and a kill switch would have stopped it in seconds. A runaway order is a violated rule, not a
small shift in a mean.
\end{solution}

\begin{solution}{exo:rs:live-experiments:6}
$-0.30 + 0.25 \times 0.883 = -0.079$, about $-0.08$. Randomised order by order, treated and control orders have the same expected share of treated
siblings, so the spillover is the same in both arms and cancels in the difference. The stock-day (all of a stock's orders on a day in one arm) contains the
\omterm{def:rs:live-experiments:ab}{interference}, and measures the rollout.
\end{solution}

\begin{solution}{exo:rs:live-experiments:7}
Median stopping days: 56 with $\tau = 0.1$, 38 with $\tau = 0.3$, 44 with $\tau = 1$. A small $\tau$ puts the alternative's mass on effects too small to
distinguish quickly; a large one spreads it over effects much larger than 0.3 and pays a factor $\sqrt{V/(V+\tau^2)}$ for them. The best $\tau$ is near the
effect one expects.
\end{solution}

\begin{solution}{exo:rs:live-experiments:8}
Daily looks at a fixed-horizon test with no effect at all reach $p < 0.05$ at some look within 23 days in 25\% of experiments, not 5\%. Either fix the
horizon and look once, or use an always-valid p-value (the mixture \omterm{def:rs:live-experiments:mde}{sequential test}), whose size holds at every look.
\end{solution}

\begin{solution}{pb:rs:live-experiments:1}
\begin{enumerate}
\item Mean 11.7, standard deviation 23.1 basis points: the pre-trade estimate, its error, the signed market move over the window, idiosyncratic noise and a
day effect.
\item Deterministic, stateless, identical on every server, independent across experiments, reproducible afterwards: no shared table to fall out of sync.
\item 2\,485 against 2\,400 planned: $z = 1.83$, p-value 0.07, no mismatch.
\item The direct saving, $-0.30$; the rollout saves $-0.08$.
\item That every server ran the new code (a second review), that alerts reached someone, and that each stage's limits and invariants capped what it could do.
\item 77 days (the formula gives 77.7).
\item The pre-trade estimate and the market's move over the scheduled window; $R^2 = 0.51$; 38 days.
\item 72 and 70 days: the pre-trade estimate explains only 9\% of the variance.
\item A spread of 0.66 and 3\,253 days for 0.3 basis point: one unit a day, and the day effect is the noise.
\item Spread 0.093 and 35 days for 0.3 basis point, like the best order-level design; it estimates the rollout's $-0.08$, which would take 503 days to detect.
\item 29\% and 33\%.
\item 46 and 62 of 4\,000.
\item Half by day 38, 4\% by day 10, 98\% by day 120.
\item At a median of 6 days.
\item \textbf{Named result.} 77 trading days at half the flow without \omterm{def:rs:live-experiments:cuped}{CUPED}, 38 with it (the formula: 77.7 and 38.2).
\item On its own orders, yes: $-0.30$ with a standard error of 0.085 in 60 days with \omterm{def:rs:live-experiments:cuped}{CUPED}. For the firm, barely: switching everything saves 0.08, not
distinguishable from zero in two years of trading.
\item Through a canary with invariants and scaled limits, then an experiment randomised by stock-day and analysed within the day with \omterm{def:rs:live-experiments:cuped}{CUPED}, a sequential rule,
guardrails on fill rate and mark-outs, and a rollback path at each stage; and ask whether a 0.08 saving is worth the change.
\item 194 orders and 8.1 days: statistics are for small shifts; rely on per-order invariants, limits sized to the canary and a kill switch for large defects.
\item The metric and guardrails, the unit, the share, the covariates, the horizon or sequential rule and $\tau$, the \omterm{def:rs:live-experiments:mde}{minimum detectable effect}, the decisions at
each outcome, logged in the \omterm{def:rs:the-research-process:log}{research log}.
\item The difference the change makes to the \omterm{def:rs:live-experiments:ab}{randomisation units} under the assignment that was run, which equals the effect of switching everything only if
no unit's outcome depends on another's assignment.
\end{enumerate}
\end{solution}

\begin{solution}{iq:rs:live-experiments:1}
\omterm{def:rs:live-experiments:paper}{Paper trading} for parity (signals, risk checks, a full calendar); a canary with verified deployment, per-order invariants, limits scaled to its share and a
kill switch; an \omterm{def:rs:live-experiments:ab}{A/B test} with a \omterm{def:rs:live-experiments:ab}{randomisation unit} that contains \omterm{def:rs:live-experiments:ab}{interference}, \omterm{def:rs:live-experiments:cuped}{CUPED} and a sequential rule; staged expansion with a tested rollback. Knight's
2012 loss came from a staged deployment whose stage check (one of eight servers) failed and whose warnings were not alerts.
\end{solution}

\begin{solution}{iq:rs:live-experiments:2}
A regression adjustment on covariates unaffected by the treatment: variance falls by $1 - R^2$ (Bing reported about half; here 0.51 with the market move).
It fails if a covariate is affected by the treatment (duration, fill rate: bias), or explains little (the pre-trade estimate alone: 9\%).
\end{solution}

\begin{solution}{iq:rs:live-experiments:3}
Each look is another chance for a false positive: daily looks give 29\% false positives in 38 days. Fix the horizon, or use always-valid p-values from a
mixture sequential probability ratio test, which bound the error over every look (Ville's inequality).
\end{solution}

\begin{solution}{iq:rs:live-experiments:4}
\omterm{def:rs:live-experiments:ab}{Interference}: the new algorithm's saving came partly at the expense of the firm's other orders in the same names (it traded earlier and left them more impact).
Order-level randomisation cancels the spillover between the arms; the rollout pays it. Randomise by stock-day (or client, or day) to measure the rollout.
\end{solution}

\begin{solution}{iq:rs:live-experiments:5}
Hash a stable identifier with the experiment's name (SHA-256, the first bytes as a fraction below the share): the same order gets the same arm on every server
without shared state, and experiments are independent. Check the split with the sample-ratio-mismatch test before reading any result, and investigate a
mismatch before trusting the comparison.
\end{solution}

\begin{solution}{iq:rs:live-experiments:6}
$n = ((z_{1-\alpha/2} + z_{1-\beta})\sigma/\Delta)^2 (1/s + 1/(1-s))$; at 5\% and 80\% the factor is 2.80. With $\sigma = 23.1$, $\Delta = 0.3$ and an equal
split: $n = (2.80 \times 23.1/0.3)^2 \times 4 \approx 186\,500$ orders, 77.7 days of 2\,400.
\end{solution}
